\subsection{Workflow: distributional improvement and decision loss diverge}
On the 61 frozen final workflow origins, the development-selected Gaussian
global model reduces macro count-fraction CRPS from 0.05010 to 0.04595, an
8.3\% improvement over independent sampling. Its mean paired difference is
$-0.00415$; the 14-day block interval is $[-0.00638,-0.00188]$, and the 28-day
interval also excludes zero. Mean 90\% coverage rises from 0.536 to 0.890,
while normalized interval width rises from 0.140 to 0.222. This illustrates a
practical failure of independent sampling: reasonable individual probabilities
can induce overly concentrated collective forecasts.

The projection breakdown localizes the effect. Population CRPS improves from
0.03421 to 0.02728 and population coverage from 0.219 to 0.869. Handler-group
CRPS improves more modestly, from 0.06599 to 0.06462, with coverage rising from
0.853 to 0.911. Equal projection weighting keeps these effects visible rather
than allowing the many handler groups to obscure the population result.

That improvement does not extend to the specified reserve loss. At under/over
ratio 9, the selected model's loss is 0.12899 versus 0.11770 for independence,
a 9.6\% increase. The paired difference is 0.01129; its 14-day block interval
is $[-0.00117,0.02129]$, while the 28-day interval is $[0.00426,0.01746]$.
The direction of the mean tradeoff is clear, but its uncertainty depends on
the block choice. We therefore report it as a held-out decision-loss tradeoff,
not a universally established cost penalty. At ratio 3, the mean loss likewise
increases, from 0.08933 to 0.09873.

Observed handler incidence provides a smaller, distinct signal. Gaussian
static factors obtain CRPS 0.04649, compared with 0.04694 after risk-stratified
incidence permutation. The difference, $-0.000454$, has a 14-day block interval
$[-0.000726,-0.000188]$. However, the corresponding reserve loss increases
from 0.11936 under permutation to 0.12219 under true incidence; the paired
interval for this increase is $[0.00119,0.00442]$. Thus, even the same
relation-incidence control can favor one model under a distributional score
and the other under a decision loss. It does not identify a causal effect of
actual handler assignment.

The richer context model does not establish a reliable improvement over the
static model: its mean CRPS difference is $-0.000188$, with a 14-day interval
$[-0.000626,0.000262]$. The common-shock static model has CRPS 0.05026 versus
0.05013 for its independent control. Merely adding a dependence model is
therefore insufficient, and the choice of family matters.

The direct historical-residual comparator achieves the lowest aggregate CRPS
among these workflow comparisons, 0.04387, while its coverage is 0.831 and its
ratio-9 loss is 0.12576. It can shift aggregate means and is not a
matched-marginal ablation. Its mean CRPS advantage over the selected Gaussian
model is not resolved by the block intervals: the Gaussian-minus-residual
difference is 0.00208, with a 14-day interval $[-0.00070,0.00536]$.
This distinction is consequential: at 14 days, the
selected individual ensemble predicts an average population resolution
fraction of 0.5252 versus the observed 0.5792. Every matched-marginal copula
inherits that mean bias. The residual baseline's advantage is evidence for
retaining simple aggregate-error controls, rather than evidence that its
separately sampled projections form a better entity-level joint law.
